% ==============================================================================
% SECTION I: INTRODUCTION
% Target Venues: IEEE S&P / ACM CCS / USENIX Security / NDSS
% Follows Academic Systems Contributions Standard (Algorithms vs Systems)
% ==============================================================================

\section{Introduction}
\label{sec:intro}

The explosive proliferation of the Internet of Things (IoT) has embedded billions of connected sensors, programmable logic controllers (PLCs), smart energy meters, and autonomous edge devices into mission-critical infrastructures~\cite{segurola2024unsupervised,dinh2020federated}. While these decentralized ecosystems unlock unprecedented operational intelligence, they also drastically expand the threat surface against sophisticated cyber attacks~\cite{sommer2010outside,sarhan2023evaluating}. In contemporary enterprise and industrial settings, edge gateways typically monitor multi-tenant administrative domains that host heterogeneous device fleets running distinct, non-interoperable application protocols (e.g., Modbus/DNP3 in power grids vs. BACnet in building automation vs. MQTT/CoAP in environmental telemetry)~\cite{ferrag2022edgeiiotset}. 

Under these operational conditions, traditional centralized Network Intrusion Detection Systems (NIDS) exhibit severe limitations. Aggregating raw packet payloads or NetFlow telemetry at a central data center incurs prohibitive wide-area network (WAN) bandwidth consumption and introduces critical transmission latency~\cite{wang2019adaptive}. More critically, pooling uncurated telemetry from multi-tenant edge gateways compromises client confidentiality and violates rigorous data sovereignty and privacy mandates, such as the General Data Protection Regulation (GDPR)~\cite{truex2019hybrid,nasr2019comprehensive}.

To circumvent these privacy and bandwidth bottlenecks, Federated Learning (FL)~\cite{mcmahan2017communication,li2020federated} has emerged as an attractive collaborative paradigm for decentralized NIDS. In a federated framework, edge gateways act as local clients that train detection models exclusively on their locally captured traffic profiles without exposing raw flows, periodically transmitting only model updates to a coordinating parameter server~\cite{rey2022federated,wang2022fedod}. However, applying FL to real-world edge IoT networks encounters an intrinsic structural barrier: \textbf{Extreme Non-IID Manifold Heterogeneity}~\cite{hsu2019measuring}. Because each gateway oversees a specialized IoT domain, the nominal traffic at gateway $\mathcal{G}_i$ is supported on a low-dimensional Riemannian sub-manifold $\mathcal{M}_i \subset \mathbb{R}^D$ that is topologically disjoint from peer manifolds in the ambient feature space:
\begin{equation}
\label{eq:intro_disjoint}
\mathcal{M}_i \cap \mathcal{M}_j \approx \emptyset \quad \forall i \neq j.
\end{equation}

Simultaneously, edge networks face a dual-spectrum threat model~\cite{mirsky2018kitsune,sarhan2023evaluating}. On one end of the spectrum lie macroscopic, high-volume Distributed Denial-of-Service (DDoS) inundations (e.g., TCP/UDP floods generated by Mirai and Bashlite botnets), which generate massive statistical feature surges pushing attack vectors deep into out-of-distribution (OOD) space. On the other end lie microscopic, stealthy multi-stage attacks (e.g., slow horizontal port scanning, asynchronous command-and-control beaconing, and lateral movement), which deliberately mimic normal marginal distributions and deviate only in subtle relational dependencies~\cite{goodge2022lunar,jin2021anemone}.

\subsection{Limitations of Existing Paradigms}
Prevailing unsupervised and one-class edge intrusion detection methodologies fail to navigate this multi-tenant Non-IID landscape effectively:
\begin{enumerate}
    \item \textbf{Reconstruction Shortcut Trap in Autoencoders}: Reconstruction-based deep architectures (e.g., Kitsune~\cite{mirsky2018kitsune}, Deep Autoencoders~\cite{sakurada2014anomaly}, and Memory AEs~\cite{gong2019memorizing}) attempt to identify anomalies via reconstruction error $\norm{x - \hat{x}}_2^2$. However, continuous deep autoencoders possess excessive expressive capacity. On subtle, low-rate multi-stage attacks, deep bottleneck layers generalize across sparse anomalous patterns---a pervasive pathology termed \emph{reconstruction shortcutting} or overgeneralization~\cite{gong2019memorizing}. Consequently, stealthy malicious flows are reconstructed with deceptively low residual error, yielding fatal false negatives.
    \item \textbf{Catastrophic Null-Space Erosion in Spectral Methods}: Analytical subspace approaches, such as the Kernel Null Foley-Sammon Transform (KNFST)~\cite{bodesheim2013kernel} and Local Orthogonal Component NFST (LOC-NFST)~\cite{nguyen2024locnfst}, project nominal data into an exact mathematical null space ($w^\top S_w w = 0, w^\top S_b w > 0$) with deterministic one-round federated aggregation. However, in continuous high-throughput streaming IoT environments ($N \gg D$), empirical covariance matrices become full-rank ($S_w \succ 0$ with probability $1$), eliminating the exact null space. Under non-stationary streaming drift and slow adversarial poisoning (e.g., ``boiling frog'' attacks), incremental updates erode the near-null space, causing decision-threshold collapse ($\tau \to 0$). Furthermore, kernelized NFST incurs $\mathcal{O}(N^2)$ memory scaling, consuming $434.8$\,MB to $958.2$\,MB RAM on edge devices.
    \item \textbf{Distance-Ranking Graph Neural Networks}: Graph-based distance-ranking architectures, epitomized by LUNAR (Goodge et al.~\cite{goodge2022lunar}), overcome these limitations by framing anomaly detection as learning local topological density transitions. By computing sorted Euclidean distance vectors $d(x) = [d_1(x), \dots, d_k(x)]^\top$ to $k$-nearest neighbors and learning a non-linear ranking function $f_\theta(d)$, LUNAR exhibits high sensitivity to localized relational anomalies while maintaining sub-millisecond inference speeds.
\end{enumerate}

\subsection{Fundamental Research Challenges}
Despite the promise of distance-ranking graph representations, our systematic mathematical and empirical analysis reveals two catastrophic failure modes when deploying distance-ranking anomaly detection across decentralized, Non-IID IoT edge networks:

\begin{figure}[t]
\centering
\includegraphics[width=\columnwidth]{confusion_matrices.png}
\caption{Empirical verification of Fed-LUNAR anomaly detection across four canonical IoT domains under severe Non-IID skew, demonstrating robust diagonal concentration and suppression of false alarms.}
\label{fig:confusion_matrices}
\end{figure}

\begin{itemize}
    \item \textbf{Challenge 1: Out-of-Distribution Distance Inversion and Monotonicity Breakdown}. Canonical distance-ranking models synthesize pseudo-negative training samples using a small, fixed perturbation radius $\epsilon = 0.1$~\cite{goodge2022lunar}, confining training distance support to a bounded interval $\mathcal{K}_{\text{train}} \subset [0.4, 1.4]^k$. When confronted with high-volume DDoS attacks (\texttt{BoTIoT}, \texttt{CICIoT2023}), attack feature magnitudes explode, pushing $k$-NN distance vectors into uncharted territory ($d_j \in [5.0, 60.0]$). Because standard Multi-Layer Perceptrons (MLPs) with LeakyReLU activations lack global monotonicity constraints, unconstrained linear layers extrapolate with inverted gradients:
    \begin{equation}
    \label{eq:intro_inversion}
    \nabla_d f_\theta(d) \cdot \frac{d}{\norm{d}_2} \le -\gamma < 0 \quad \text{as } \norm{d}_2 \to \infty,
    \end{equation}
    driving predicted anomaly logits toward $-\infty$ ($\sigma(f_\theta) \to 0.0000$). Consequently, massive volumetric attacks are classified as more nominal than legitimate traffic, causing empirical AUC-ROC to catastrophically collapse to **0.15\%** on \texttt{BoTIoT} and **4.58\%** on \texttt{CICIoT2023}.
    \item \textbf{Challenge 2: Cross-Manifold Negative Gradient Cancellation in Non-IID FL}. In a decentralized multi-tenant deployment, edge gateways independently generate pseudo-negative vectors to delineate their local benign boundaries. Because client manifolds are mutually disjoint ($\mathcal{M}_A \cap \mathcal{M}_B \approx \emptyset$), uncoordinated isotropic perturbations generated by Gateway $\mathcal{G}_A$ frequently fall directly inside the legitimate traffic manifold of peer Gateway $\mathcal{G}_B$ ($\tilde{x}_A \in \mathcal{M}_B$). Client $\mathcal{G}_A$ computes loss gradients pushing $\sigma(f_\theta(\tilde{x}_A)) \to 1$ (malicious), whereas Client $\mathcal{G}_B$ computes gradients pushing $\sigma(f_\theta(x_B)) \to 0$ (benign). We prove that their inner product is strictly negative:
    \begin{equation}
    \label{eq:intro_cancellation}
    \inner{\nabla_\theta \mathcal{L}_A(\tilde{x}_A)}{\nabla_\theta \mathcal{L}_B(x_B)} < 0.
    \end{equation}
    Under standard Federated Averaging (FedAvg)~\cite{mcmahan2017communication}, opposing client gradients annihilate each other ($\norm{\nabla_\theta \mathcal{L}_{\text{global}}} \to 0$). Empirical logs demonstrate that pairwise gradient conflicts occur in up to **70.0\%** of federated rounds under Dirichlet Non-IID skew ($\alpha = 0.5$), severely depressing Macro F1 to **57.22\%**.
\end{itemize}

\subsection{Summary of Contributions}
Following rigorous standards for systems and machine learning security research, we delineate our scientific contributions into algorithmic foundations and systems/hardware-software co-design:

\subsubsection{Algorithmic Contributions}
\begin{enumerate}
    \item \textbf{Theoretical Proof of OOD Distance Inversion and Monotonicity Recovery (Theorem 1)}: We mathematically characterize the geometry of distance-ranking extrapolation breakdown in piecewise-affine LeakyReLU networks, formally proving the mechanism of AUC collapse ($0.15\%$). We propose \emph{Multi-Scale Subspace Perturbation (MSSP)}, which injects synthetic negatives across a geometrically spaced radial spectrum ($\mathcal{S}_{\text{scales}} = \{0.2, 0.5, 1.5, 3.0, 6.0\}$) spanning the metric diameter up to $R_{\max}$, and prove that MSSP guarantees strictly positive rank monotonicity:
    \begin{equation}
    \frac{\partial f_\theta(d)}{\partial d_j} > 0 \quad \forall j \in \{1, \dots, k\}, \quad \forall d \in [0, R_{\max}]^k.
    \end{equation}
    \item \textbf{Theoretical Proof of Gradient Cancellation and Purging Invariance (Theorem 2 \& Lemma 2.1)}: We formally prove the mathematical dynamics of adversarial negative gradient cancellation arising from uncoordinated cross-manifold intrusion. To resolve this without transmitting private raw flows, we design \emph{Federated Subspace Density Sketches (FSDS)} and \emph{Cross-Manifold Negative Purging (CMNP)}. Each client broadcasts a compact quadruple $\mathcal{S}_i = (\mu_i, \Lambda_i, U_i, r_{i,\max})$ occupying $<5$\,KB. Peer clients filter candidate negatives against peer sketches via a joint null-space proximity and Mahalanobis containment test, proving that invasive candidates are purged with confidence $\ge 99\%$.
    \item \textbf{Conflict-Averse Gradient Alignment (Lemma 2.2)}: To eliminate residual optimization tensions arising from scale disparities across heterogeneous edge tenants, we develop \emph{Distance-Ranking Orthogonal Gradient Alignment (DROGA)}. DROGA normalizes client updates onto the unit hypersphere and solves a dual simplex quadratic program (DR-CAGrad QP) per communication round at the central server, proving that the aligned update $g_{\text{aligned}}$ guarantees non-negative projection with all participating clients ($\inner{g_{\text{aligned}}}{g_i} \ge 0$) and drives monotonic Pareto-objective descent under $L$-smoothness.
\end{enumerate}

\subsubsection{Systems and Hardware-Software Co-Design Contributions}
\begin{enumerate}
    \item \textbf{Sub-Millisecond Line-Rate Edge Streaming Engine}: We implement an end-to-end edge-native inference pipeline optimized for commodity single-board computing gateways (e.g., ARM Cortex-A72 / Raspberry Pi 4). By decoupling lightweight $k$-NN distance retrieval from asynchronous federated aggregation, Fed-LUNAR achieves deterministic per-sample inference latency of \textbf{0.0008\,ms to 0.0020\,ms} (throughput $>500{,}000$ packets/second). This enables line-rate inspection on $1$\,Gbps to $10$\,Gbps industrial backbones without queue congestion or packet dropping.
    \item \textbf{Ultra-Low Memory Footprint}: Unlike spectral and kernelized null-space methods that require $\mathcal{O}(N^2)$ Gram matrices and consume up to $958.2$\,MB RAM, Fed-LUNAR maintains an active runtime memory working set of only \textbf{49.6\,MB to 67.1\,MB RAM}---a quantifiable reduction of \textbf{85\% to 93\%}, ensuring deployment feasibility on memory-constrained edge hardware.
    \item \textbf{Comprehensive Multi-Dataset Physical Benchmark}: We benchmark Fed-LUNAR across four canonical IoT datasets (\texttt{BoTIoT}, \texttt{EdgeIIoTset}, \texttt{CICIoT2023}, and high-dimensional 115-feature \texttt{N\_BaIoT}) executed on dedicated NVIDIA RTX 5090 hardware. Across 32 full benchmark runs under Dirichlet non-identical distribution parameters $\alpha \in \{0.1, 0.5, 1.0, 5.0\}$, Fed-LUNAR consistently maintains high F1-Scores ($93.14\%$ on \texttt{BoTIoT}, $97.54\%$ on \texttt{N\_BaIoT}), remaining remarkably resilient even at extreme statistical heterogeneity ($\alpha=0.1$). Ablation analysis confirms that removing CMNP precipitates an immediate $26.23\%$ collapse in F1-Score, empirically validating our theoretical formulations.
\end{enumerate}

\subsection{Paper Organization}
The remainder of this paper is structured as follows. Section~\ref{sec:threat_model} establishes the multi-tenant IoT edge system model, formalizes the threat model, and articulates the defensible research gap via the Anomaly Detection Trilemma. Section~\ref{sec:formulation} formulates the mathematical mechanics of graph distance-ranking and gradient conflict. Section~\ref{sec:methodology} presents the complete Fed-LUNAR methodology, detailing MSSP, FSDS, CMNP, and DROGA. Section~\ref{sec:proofs} provides full step-by-step mathematical proofs of Theorems~1 and~2, as well as Lemmas~2.1 and~2.2. Section~\ref{sec:experiments} presents multi-dataset empirical benchmarks, Dirichlet sweeps, ablation studies, and edge hardware profiling from server execution logs. Section~\ref{sec:related} positions Fed-LUNAR across a five-paradigm taxonomy with an exhaustive comparison matrix. Finally, Section~\ref{sec:conclusion} concludes the paper with directions for future edge research.
